\documentclass[sigconf]{acmart}

\settopmatter{printfolios=true}
\copyrightyear{2026}
\acmYear{2026}
\acmConference[DS Research Project]{Data Science Research Project}{March 2026}{Sri Lanka}
\acmBooktitle{Data Science Research Project}
\acmDOI{10.1145/1122445.1122456}
\acmISBN{978-1-4503-XXXX-X/18/06}

\PassOptionsToPackage{hyphens}{url}
\usepackage{hyperref}
\sloppy

\begin{document}

\title{A Data-Driven Analysis of Migration and Remittance: Insights for Related Decision-Making in Sri Lanka}

\author{Ginige D.N.}
\email{dinurag.23@cse.mrt.ac.lk}
\affiliation{%
  \institution{University of Moratuwa}
  \streetaddress{Department of Computer Science \& Engineering}
  \city{Moratuwa}
  \country{Sri Lanka}
}

\author{Lakshan H.M.K.}
\email{kalanal.23@cse.mrt.ac.lk}
\affiliation{%
  \institution{University of Moratuwa}
  \streetaddress{Department of Computer Science \& Engineering}
  \city{Moratuwa}
  \country{Sri Lanka}
}

\author{Fernando K.N.D.}
\email{dhinanjayaf.23@cse.mrt.ac.lk}
\affiliation{%
  \institution{University of Moratuwa}
  \streetaddress{Department of Computer Science \& Engineering}
  \city{Moratuwa}
  \country{Sri Lanka}
}

\author{Gurusinghe C.R.}
\email{chanupag.23@cse.mrt.ac.lk}
\affiliation{%
  \institution{University of Moratuwa}
  \streetaddress{Department of Computer Science \& Engineering}
  \city{Moratuwa}
  \country{Sri Lanka}
}

\author{Pahanga G.H.L.}
\email{lasanap.23@cse.mrt.ac.lk}
\affiliation{%
  \institution{University of Moratuwa}
  \streetaddress{Department of Computer Science \& Engineering}
  \city{Moratuwa}
  \country{Sri Lanka}
}

\begin{abstract}
This research project aims to conduct a comprehensive data-driven analysis of migration patterns and remittance flows in Sri Lanka. By leveraging data science methodologies and machine learning techniques, we seek to uncover hidden patterns, correlations, and predictive insights that can inform policy decisions and economic planning. The study utilizes multiple data sources including Central Bank of Sri Lanka records, World Bank open data, Federal Reserve economic data, and ILO migration statistics, spanning approximately 25 years. Through exploratory data analysis (EDA), time series analysis, and predictive modeling, we expect to identify key demographic and economic factors influencing emigration and remittance trends, ultimately contributing to evidence-based decision making for sustainable national development.
\end{abstract}

\keywords{Migration, Remittances, Data Science, Machine Learning, Time Series Analysis, Vector Autoregression, Sri Lanka}

\maketitle

\section{Problem Motivation}

In Sri Lanka, emigration is a phenomenon with varying trends for the last few decades, while a significant count of the population with working capability move outwards. Consequently, this movement generates substantial inward remittances, which help in a country's economic stability. Remittances have direct contributions to household income, poverty reduction, and foreign exchange reserves, and indirectly they give rise to complex dependencies that influence the country's domestic development. Despite their importance, the underlying demographic and economic factors driving migration and remittance flows remain crucial to study by data-science approach.

The study of emigration and remittance flows is essential because remittances normally contribute to 6\% - 9\% of GDP of Sri Lanka\cite{remittance-gdp}. Remittances are not just private transfers but they act as a stabilizing mechanism for foreign reserves, household consumption, and poverty alleviation while being shaped by external shocks, such as fluctuations in global labor demand or exchange rates. A data-science approach will be able to identify hidden demographic and economic patterns, offering predictive insights of future emigration and remittance trends. Such evidence-based analysis is critical to design sustainable emigration strategies to secure the country's domestic human resources, manage financial prospects, and for communities to reduce vulnerability to economic uncertainty. Subsequently, the research directly contributes to informed decision-making and strengthens Sri Lanka's capacity for inclusive growth, sustainability and promising inflows.

The beneficiaries of such analysis are multiple. Policymakers can gain actionable knowledge, financial institutions can anticipate remittance flows and develop tailored products, families and communities can benefit from reliable support programs and reduce their vulnerability to economic shocks and researchers acquire new evidence to advance related studies based on Sri Lanka. Exploring hidden trends and building predictive models is the main aim of this project which generates insights that strengthen decision-making and contribute to sustainable national development.

Existing literature has examined Sri Lanka's remittance trends from macroeconomic and labour migration perspectives\cite{ilo-migration}. However, these studies rely on descriptive and econometric methods applied to limited time windows, and few leverage multi-source datasets or modern machine learning approaches. This project fills that gap by integrating diverse data sources spanning 25 years with a comprehensive analytical pipeline combining time series modeling, VAR-based causal analysis, and machine learning.

\section{Identified Data Sources}

\textbf{Existing Datasets:} The main datasets identified include CBSL monthly and quarterly workers' remittances from 2009 to present with country and skill-level breakdowns \cite{worldbank-remittances}, LKR/USD exchange rates \cite{cbsl-exchangerate}, and the Colombo Consumer Price Index (CCPI) \cite{ccpi}. SLBFE annual statistical handbooks provide departures by skill category from 1986 onwards, combining the legacy system (1986--2007) and the current website (2008--2024) \cite{slbfe}. Brent Crude oil prices are sourced from the EIA \cite{eia-oil}. World Bank open data covers personal remittances (current USD and \% of GDP) from 1975--2024 via KNOMAD, and GDP per capita at PPP \cite{worldbank-ppp}. FRED provides remittance inflows to GDP annually from 1975--2020 \cite{fred-remittances}. Trading Economics supplies aggregated monthly CBSL inflow data (2009--2025) with forecasts, and the ILO supplies standardized migration and remittance statistics \cite{ilo-migration}.

\textbf{APIs and Integrated Datasets:} The World Bank and FRED datasets are accessible via open APIs, enabling programmatic retrieval for reproducibility. Since some data are expressed as graphs we are planning to form integrated datasets by combining those graph related data and other types of datasets. The integrated dataset will merge remittance inflow series with macroeconomic indicators (USD/LKR exchange rates, oil prices, GDP growth, CPI, PPP) and migration flow statistics, resolving inconsistencies in temporal granularity and units across sources.

\textbf{Expected Size and Type of Data:} With related to remittances we expect 25 years past monthly remittance data which are numerical. The data related to migration are also needed covering at least the past 25 years which also mostly include numerical data. In total, the integrated dataset is expected to contain several thousand observations across approximately 10--15 numerical features, with no individual-level data.

\textbf{Data Documentation:} Key limitations include that CBSL granular records only begin from 2009, requiring World Bank and FRED sources to extend coverage to 1975. SLBFE revised its classification to only reflect skilled and low-skilled categories, discontinuing finer breakdowns, affecting longitudinal comparability. Some variables are available only for shorter windows or in graphical form, introducing potential extraction error. All datasets are publicly available and ethically sourced. No individual-level data will be used, ensuring compliance with data privacy standards.

\section{Proposed Methodology}

This study will employ a quantitative, data-driven methodology, leveraging techniques from data science and machine learning to analyze data related to Sri Lankan emigration, remittances and underlying features.

\subsection{Data Analysis Techniques}

The analysis will proceed in two main phases: exploratory data analysis (EDA) and predictive modeling. Initially, Exploratory Data Analysis (EDA) will be conducted to summarize the dataset's main characteristics, identify initial patterns, detect outliers, and test underlying assumptions. This will involve statistical analysis and data visualization to understand distributions and correlations between key variables such as demographics (age, education, district of origin), economic indicators (GDP growth, unemployment rates, exchange rates), and remittance volumes. Subsequently, Time Series Analysis will be employed to decompose remittance and migration data into trend, seasonal, and residual components. This will help in understanding the cyclical nature of flows and their correlation with external economic shocks and policy changes. Stationarity of all series will be verified using ADF and KPSS tests prior to modeling, with differencing applied where necessary. Both exploratory visualizations (correlation heatmaps, time series plots) and explanatory visualizations (trend dashboards, impulse response plots, cluster profiles) will be produced to guide analysis and communicate findings to policymakers.

\subsection{Machine Learning Approaches}

To move beyond descriptive analysis and generate predictive insights, several machine learning techniques will be explored:

\begin{itemize}
    \item \textbf{Regression Models:} To predict remittance volumes based on economic indicators (USD/LKR exchange rates, GDP growth, oil prices, CPI, PPP) and demographic factors, enabling forecasting of future trends. Methods such as Ridge, Lasso, and Gradient Boosting will be evaluated using RMSE and MAE on held-out test sets.

    \item \textbf{Clustering Algorithms:} To identify distinct patterns in migrant profiles and remittance behaviors, grouping similar characteristics for targeted policy interventions. K-means and hierarchical clustering will be applied across features such as destination country, skill level, and remittance volume.

    \item \textbf{Classification Models:} To categorize migration patterns and predict likelihood of emigration based on demographic and economic variables, such as periods of high versus low remittance inflows relative to historical baselines.

    \item \textbf{Vector Autoregression (VAR):} To model the interdependencies between remittance flows and key macroeconomic variables such as USD/LKR exchange rates, oil prices, GDP growth, CPI, and PPP simultaneously. Since remittances and these indicators are likely to influence each other over time, VAR captures these bidirectional relationships that standard regression cannot. Granger causality tests will be used to formally assess which variables drive remittance flows, and impulse response functions (IRF) will be derived to analyze how shocks---such as the 2022 economic crisis---propagate across the system. If long-run cointegration is found among variables, a Vector Error Correction Model (VECM) will be considered as an extension.
\end{itemize}

All models will be evaluated using appropriate metrics---RMSE and MAE for regression, silhouette scores for clustering, and accuracy/F1 for classification---with cross-validation applied to prevent overfitting.

\section{Expected Outputs}

\begin{itemize}
    \item \textbf{Integrated Dataset:} A cleaned, unified multi-source dataset spanning 25 years, combining CBSL, World Bank, FRED, ILO, and Trading Economics data, reusable for future researchers.

    \item \textbf{Hidden Patterns:} Correlations between macroeconomic factors and remittance inflow fluctuations, and demographic shifts in migrant profiles.

    \item \textbf{Predictive Insights:} A predictive model to forecast remittance trends as a result of factor changes.

    \item \textbf{Visualizations:} Exploratory and explanatory visualizations including trend plots, cluster profiles, and impulse response charts for academic and policy audiences.
\end{itemize}

\vspace{12pt}
\bibliographystyle{ACM-Reference-Format}
\bibliography{references}

\end{document}